% =============================================================================
% 3 The capacity model. Ben, 2026-10-18 (review notes 8, 9, 10): the model is
% where the decision is made; admission (4.2) only executes it. Define every
% quantity, say how F and c are obtained, and present the operating point as an
% APPLICATION of the model to a measured T0, not as an estimate. Context cancels
% in the WORKLOAD view (3.4) only; it does not cancel in admission (4.2).
% =============================================================================

\section{The Capacity Model}\label{sec:model}

The capacity model decides, for one model on one node, whether riders pay, how
many of them pay, and how many bytes they may move per step. It is built from a
profile taken once per model and node, plus one quantity, the baseline step
time $T_0$, that is measured on the serving node itself. Admission
(\S\ref{sec:admission}) then executes its decision every step.

\begin{table}[htb]
\centering\small\setlength{\tabcolsep}{3.5pt}
\caption{\textbf{The capacity model's inputs.} The first group describes the
model and the node; the second is the feature's profile (\S\ref{sec:profile}).}
\label{tab:symbols}
\begin{tabular}{@{}l >{\raggedright\arraybackslash}p{3.9cm} >{\raggedright\arraybackslash}p{2.2cm}@{}}
\toprule
Symbol & Meaning & Obtained by \\
\midrule
$L$, $\btok$ & layers; KV bytes per token, $2L\,h_{\mathrm{kv}}\,d\,s$ & model config \\
$\Mkv$ & HBM left for KV after weights and activations & engine's memory profile \\
$B$ & residents that fit: $\sum_{\mathrm{res}} S_i\btok \le \Mkv$ & from $\Mkv$ \\
$T_0$ & step time with $B$ residents, no riders & measured on the node \\
$\Xlink$ & host-to-device bandwidth the copy stream achieves & copy benchmark \\
\midrule
$F$ & fixed cost per step once any rider exists & profile fit \\
$c$ & exposed cost of each additional rider & profile fit \\
\bottomrule
\end{tabular}
\end{table}

\subsection{Profiling the Overhead: $F$ and $c$}\label{sec:profile}

Riders slow the step. The profile measures by how much, as a function of the
rider count $r$. It holds a full resident cohort fixed, adds riders one at a
time, and records the step time $T(r)$ over decode-only steps; the overhead is
$\Delta T(r) = T(r) - T_0$. Below the knee defined in \S\ref{sec:pays}, the
overhead is affine:
\begin{equation}
  \Delta T(r) \;=\; F + c\,r .
  \label{eq:law}
\end{equation}
The two coefficients are different kinds of work. $F$ is paid once per step as
soon as one rider exists, independent of how many: per layer, the second
attention call over the rider rows and the write of the new token's $K,V$ to the
host copy (\S\ref{sec:mm}). $c$ is what each rider adds that cannot overlap
with compute. It has four parts. The \emph{mechanism} part, $c_{\mathrm{m}}$, is
the rider's per-layer copy launches and its rows in that attention call; it
depends neither on context nor on link speed. The second part is the rider's
attention itself, which reads its $S$ keys and values from the device slot, the
same HBM read a resident's attention performs. The third is the price of
landing those bytes in HBM. Below the knee the transfer's \emph{time} hides
behind the previous layer, but its \emph{bytes} are written into HBM while the
step runs, and a decode step is bound by HBM bandwidth: every byte the copy
writes is bandwidth the step's own reads lose. The fourth is the one part of
the transfer that cannot hide: the first layer has no previous layer to hide
behind, so $1/L$ of the rider's bytes cross the link in the open. So
\begin{equation}
\begin{aligned}
  c(S) &\;=\; c_{\mathrm{m}} + \frac{S\btok}{\BWkv} + \kappa\,\frac{S\btok}{\mathit{BW}}
           + \frac{S\btok}{L\,\Xlink},\\
  T_0 &\;=\; \frac{W}{\BWw} + \sum_{\mathrm{res}} \frac{S_i\btok}{\BWkv},
\end{aligned}
  \label{eq:cS}
\end{equation}
where $\BWw$ and $\BWkv$ are the HBM bandwidths that weight and KV reads
achieve, $\mathit{BW}$ is the HBM peak, and $\kappa$ is the step time lost per
unit of HBM time the copy's writes occupy. The baseline step pays the
attention term for each resident. The contention term does not depend on the
link's speed: a faster link writes the same bytes in less time, at a
proportionally higher rate. We measure $\kappa = 1.06$ on A100 for bulk copies
(\S\ref{sec:ablation}). DAK observes the same on GH200: background copies
degrade local memory performance by up to 10\%, and the degradation scales
with the ratio of link to HBM bandwidth~\cite{dak}, which is what the term
expresses.
% src: dak s2.3: "This degradation scales with the interconnect-to-local
%      bandwidth ratio. Background copies can degrade local memory performance
%      by up to 10% on a GH200". E1 (job 2191314): 0.055% per GB/s of copy
%      traffic x 1935 GB/s = 1.06. (BOOST's GH200 numbers deliberately not
%      cited, Ben 2026-10-21.)

Three rules make the fit trustworthy. First, only points well below the knee are
kept, those whose transfer time is at most a tenth of $T_0$. A run that crosses
the knee adds exposed transfer at high $r$, which steepens the line and lowers
its intercept, overstating $c$ and understating $F$ together. Second, the profile
is repeated at several depths, and both coefficients are expressed per layer,
$F = f_L L$ and $c = c_0 + c_L L$, so one profile carries to a model of
different depth. Third, it is repeated at several contexts with the
transfer's bytes removed, which separates $c_{\mathrm{m}}$ from the attention
term of \cref{eq:cS}; $F$ does not depend on context. A single run past the knee gives the third coefficient: there the step is
the link time plus $\delta$ (\cref{eq:step}), so $\delta$ is read off as their
difference. The contention
coefficient $\kappa$ is measured directly: the
step's weight GEMMs are timed while a second stream copies into HBM at a
controlled rate (\S\ref{sec:ablation}).
% src: OVERHEAD_CHARACTERISATION.md s0a (M13): retained points load <= 0.1;
%      "fitting across the knee ... raises the slope and lowers the intercept";
%      F/c decomposition s0a M14 (mirror write -> F, streaming -> c).

\subsection{When Riders Pay}\label{sec:pays}

With $B$ residents, a step delivers $B$ tokens in $T_0$, so one token is worth
$T_0/B$ of step time. Adding $r$ riders multiplies the tokens per step by
$1+r/B$ and the step time by $1+(F+c\,r)/T_0$. The throughput gain is
\begin{equation}
  G(r) \;=\; \frac{1 + r/B}{1 + (F + c\,r)/T_0},
  \label{eq:gain}
\end{equation}
which exceeds one exactly when
\begin{equation}
  r\left(\frac{T_0}{B}-c\right) \;>\; F .
  \label{eq:pays}
\end{equation}
Three conditions follow, each of a different kind.
\begin{itemize}\itemsep1pt
\item \textbf{The gate}, $c < T_0/B$. If one rider costs more step time than
  one token is worth, no number of riders pays. The gate depends on neither $r$
  nor $F$, and it is a property of the implementation: the building blocks of
  \S\ref{sec:mm} exist to pass it.
\item \textbf{Break-even}, $r \ge \rmin$. Riders must be numerous enough to
  repay $F$. The smallest such count is
  $\rmin = \lfloor F/(T_0/B - c)\rfloor + 1$.
\item \textbf{The knee}, where the link stops hiding. With $\Xlink$ the
  host-to-device bandwidth the copy stream achieves, the riders' transfer
  takes $\sum_i S_i\btok/\Xlink$. It hides behind the step, and the step is
  not $T_0$: every admitted rider lengthens it, by $F$ for the first and $c$
  for each, so each rider widens the window the next one's transfer hides in.
  Not all of the step is window. A part of it, $\delta$, has no copy in
  flight, and the profile measures it alongside $F$ and $c$
  (\S\ref{sec:profile}). The riders hide while
  $\sum_i S_i\btok/\Xlink \le T_0 + F + c\,r - \delta$; past that point the
  link sets the step, so
  \begin{equation}
    T(r) \;=\; \max\Bigl(T_0 + F + c\,r,\;\; \sum_i \frac{S_i\btok}{\Xlink} + \delta\Bigr).
    \label{eq:step}
  \end{equation}
  For riders of equal context $S$ the knee is
  $\rstar = (T_0 + F - \delta)/(S\btok/\Xlink - c)$.
\end{itemize}
The throughput gain of \cref{eq:gain} is the general
$G(r) = (1+r/B)\,T_0/T(r)$ below the knee. When the gate holds, $G$ increases
with $r$ up to the knee (its derivative has the sign of $T_0 + F - cB$, which
the gate makes positive). Past the knee every rider adds its whole transfer and
one token, and $G$ falls as long as $\delta < B\,S\btok/\Xlink$, which holds
at every point we evaluate. The best count is therefore one of the two whole
numbers around the knee: the last rider that hides, or the first whose
transfer is partly exposed but still adds more than it costs. The decision is
\begin{equation}
  r \;=\;
  \begin{cases}
    \arg\max_{\,r \in \{\lfloor\rstar\rfloor,\, \lceil\rstar\rceil\}} G(r)
      & \text{if } c < T_0/B \text{ and } r \ge \rmin,\\
    0 & \text{otherwise.}
  \end{cases}
  \label{eq:decision}
\end{equation}
With mixed contexts there is no single count; admission applies the same test
one request at a time, adding a rider while \cref{eq:step} predicts that it
raises throughput (\S\ref{sec:admission}).
% src: Ben 2026-10-22 ("apply the 1.37x rule across the paper"). M22c (job
%      2190742): past the knee step = link*r + 6.8 ms (r = 5), 7.5 (r = 6);
%      eq. step reproduces 88.4/96.7/94.9 tok/s at r = 4/5/6 (measured 88.4,
%      96.7, 94.8; F, c and delta from the same runs) and picks r = 5, the
%      measured best. paper/model/knee_rule.py, gain_curve.py rule="max".

\subsection{What Context Does and Does Not Change}\label{sec:context}

For a workload of one context $S$, both tiers fill in proportion to $1/S$: HBM
holds $B = \Mkv/(S\btok)$ residents, and within the residents' step the link
moves the KV of $r_0 = T_0\Xlink/(S\btok)$ riders. Their ratio,
\begin{equation}
  \frac{r_0}{B} \;=\; \frac{T_0\,\Xlink}{\Mkv},
  \label{eq:ratio}
\end{equation}
is, to first order, the throughput riders add, and it does not depend on $S$;
riders' own step time and the tail $\delta$ move the knee of \cref{eq:step}
around $r_0$. Nor does $T_0$ of a
full cohort: it reads the weights and $\Mkv$ bytes of KV whatever the context.
A longer context
fits fewer residents \emph{and} fewer riders, and leaves the ratio unchanged:
what riders can add is set by the hardware, by how many bytes the link moves
in one step relative to how many bytes of KV HBM holds. What riders \emph{cost}
does depend on context, mildly and favourably: $r_0$ riders read
$r_0 S\btok = T_0\Xlink$ bytes of attention, and write as many into HBM,
whatever $S$ is, while their mechanism cost $r_0 c_{\mathrm{m}}$ falls as
$1/S$. Context also matters
in two places. When choosing a workload, the rider count is an integer
(\S\ref{sec:opoint}). When admitting from a queue of mixed contexts, it decides
\emph{which} requests should ride (\S\ref{sec:admission}).

\subsection{Applying the Model: the Operating Point}\label{sec:opoint}

To run a model on a node, the system first serves the largest resident cohort
that fits, with no riders, and measures its step time; that is $T_0$. With
$\Xlink$ and the profile, \cref{eq:gain,eq:decision} then give the gain at any
context, and admission re-applies them every step against a $T_0$ it keeps
measuring (\S\ref{sec:admission}).

% Section 3.5: gain against context at the target (Ben, 2026-10-20: "show gain
% vs context as a curve in the paper"). Data printed by paper/model/gain_curve.py:
% the M22-corrected model, c(S) = c_mech + S b_tok/BW_kv and T0 including the
% residents' KV reads; knee at T0*X. Qwen2.5-72B fp8 weights, 16-bit KV, GH200-96,
% X = 419 GB/s. Integer: whole residents and riders. Smooth: fractional B and r.
% 2026-10-21: c also carries the HBM-contention term KAPPA*S*b_tok/BW_hbm (E1,
% KAPPA = 1.06, bulk copies); smoothnc is the envelope without it.
% Below 2 residents (S > 24.4k) the configuration is degenerate and not drawn.
\begin{figure}[t]
\centering
\begin{tikzpicture}
\begin{axis}[
  width=0.95\columnwidth, height=4.8cm,
  xlabel={context $S$ (k tokens)},
  ylabel={throughput gain $G$},
  xmin=2, xmax=24, ymin=1.0, ymax=1.62,
  xtick={2,4,8,12,16,20,24},
  ytick={1.0,1.1,1.2,1.3,1.4,1.5,1.6},
  yticklabel style={/pgf/number format/.cd,fixed,fixed zerofill,precision=1},
  tick label style={font=\sffamily\scriptsize},
  label style={font=\sffamily\scriptsize},
  axis line style={black!60}, tick style={black!60},
  ymajorgrids, grid style={black!10},
  legend style={font=\sffamily\scriptsize, draw=none, fill=none,
                at={(0.02,0.98)}, anchor=north west, fill=white, fill opacity=0.85, text opacity=1},
  legend cell align=left,
]
\pgfplotstableread{
S  integer  smooth  smoothnc
    2.0 1.0158 1.0103 1.0764
    2.1 1.0240 1.0210 1.0886
    2.2 1.0324 1.0310 1.0999
    2.3 1.0411 1.0403 1.1105
    2.4 1.0499 1.0489 1.1203
    2.5 1.0668 1.0570 1.1296
    2.6 1.0771 1.0646 1.1382
    2.7 1.0652 1.0717 1.1463
    2.8 1.0846 1.0784 1.1540
    2.9 1.0956 1.0847 1.1612
    3.0 1.0926 1.0906 1.1680
    3.1 1.1039 1.0962 1.1745
    3.2 1.1009 1.1016 1.1806
    3.3 1.1270 1.1066 1.1864
    3.4 1.1096 1.1114 1.1919
    3.5 1.1381 1.1160 1.1972
    3.6 1.1351 1.1203 1.2022
    3.7 1.1159 1.1245 1.2069
    3.8 1.1470 1.1284 1.2115
    3.9 1.1442 1.1322 1.2158
    4.0 1.1414 1.1358 1.2200
    4.1 1.1567 1.1392 1.2240
    4.2 1.1540 1.1425 1.2278
    4.3 1.1514 1.1457 1.2314
    4.4 1.1488 1.1487 1.2350
    4.5 1.1644 1.1517 1.2383
    4.6 1.1620 1.1545 1.2416
    4.7 1.1595 1.1572 1.2447
    4.8 1.1571 1.1598 1.2477
    4.9 1.2031 1.1623 1.2506
    5.0 1.1705 1.1647 1.2534
    5.1 1.1683 1.1670 1.2561
    5.2 1.1661 1.1693 1.2587
    5.3 1.1639 1.1715 1.2613
    5.4 1.1617 1.1736 1.2637
    5.5 1.2134 1.1756 1.2660
    5.6 1.2110 1.1776 1.2683
    5.7 1.1729 1.1795 1.2705
    5.8 1.1710 1.1813 1.2727
    5.9 1.1690 1.1831 1.2747
    6.0 1.1671 1.1848 1.2768
    6.1 1.1652 1.1865 1.2787
    6.2 1.2240 1.1881 1.2806
    6.3 1.2218 1.1897 1.2824
    6.4 1.2197 1.1913 1.2842
    6.5 1.2176 1.1928 1.2860
    6.6 1.2155 1.1942 1.2877
    6.7 1.1705 1.1956 1.2893
    6.8 1.1688 1.1970 1.2909
    6.9 1.1671 1.1983 1.2924
    7.0 1.2349 1.1996 1.2940
    7.1 1.2330 1.2009 1.2954
    7.2 1.2311 1.2021 1.2969
    7.3 1.2293 1.2033 1.2983
    7.4 1.2274 1.2045 1.2996
    7.5 1.2256 1.2057 1.3010
    7.6 1.2238 1.2068 1.3023
    7.7 1.2220 1.2079 1.3035
    7.8 1.2202 1.2089 1.3048
    7.9 1.2184 1.2100 1.3060
    8.0 1.2166 1.2110 1.3072
    8.1 1.1622 1.2120 1.3083
    8.2 1.2407 1.2129 1.3095
    8.3 1.2392 1.2139 1.3106
    8.4 1.2376 1.2148 1.3116
    8.5 1.2360 1.2157 1.3127
    8.6 1.2344 1.2166 1.3137
    8.7 1.2329 1.2175 1.3147
    8.8 1.2314 1.2183 1.3157
    8.9 1.2298 1.2191 1.3167
    9.0 1.2283 1.2200 1.3176
    9.1 1.2268 1.2207 1.3186
    9.2 1.2253 1.2215 1.3195
    9.3 1.2238 1.2223 1.3204
    9.4 1.2223 1.2230 1.3212
    9.5 1.2209 1.2238 1.3221
    9.6 1.2194 1.2245 1.3229
    9.7 1.2180 1.2252 1.3238
    9.8 1.2416 1.2259 1.3246
    9.9 1.2403 1.2266 1.3254
   10.0 1.2391 1.2272 1.3261
   10.1 1.2378 1.2279 1.3269
   10.2 1.2366 1.2285 1.3276
   10.3 1.2354 1.2291 1.3284
   10.4 1.2341 1.2298 1.3291
   10.5 1.2329 1.2304 1.3298
   10.6 1.2317 1.2310 1.3305
   10.7 1.2305 1.2316 1.3312
   10.8 1.2293 1.2321 1.3319
   10.9 1.2281 1.2327 1.3325
   11.0 1.2269 1.2332 1.3332
   11.1 1.2257 1.2338 1.3338
   11.2 1.2246 1.2343 1.3344
   11.3 1.2234 1.2349 1.3350
   11.4 1.2222 1.2354 1.3356
   11.5 1.2211 1.2359 1.3362
   11.6 1.2199 1.2364 1.3368
   11.7 1.2188 1.2369 1.3374
   11.8 1.2176 1.2374 1.3380
   11.9 1.2165 1.2378 1.3385
   12.0 1.2154 1.2383 1.3391
   12.1 1.2142 1.2388 1.3396
   12.2 1.2131 1.2392 1.3401
   12.3 1.3632 1.2397 1.3407
   12.4 1.3618 1.2401 1.3412
   12.5 1.3604 1.2405 1.3417
   12.6 1.3589 1.2410 1.3422
   12.7 1.3575 1.2414 1.3427
   12.8 1.3561 1.2418 1.3432
   12.9 1.3547 1.2422 1.3437
   13.0 1.2209 1.2426 1.3441
   13.1 1.2200 1.2430 1.3446
   13.2 1.2192 1.2434 1.3450
   13.3 1.2183 1.2438 1.3455
   13.4 1.2175 1.2442 1.3459
   13.5 1.2166 1.2445 1.3464
   13.6 1.2158 1.2449 1.3468
   13.7 1.2150 1.2453 1.3472
   13.8 1.2141 1.2456 1.3477
   13.9 1.2133 1.2460 1.3481
   14.0 1.2125 1.2463 1.3485
   14.1 1.2117 1.2467 1.3489
   14.2 1.2109 1.2470 1.3493
   14.3 1.2100 1.2474 1.3497
   14.4 1.2092 1.2477 1.3501
   14.5 1.2084 1.2480 1.3504
   14.6 1.2076 1.2483 1.3508
   14.7 1.2068 1.2487 1.3512
   14.8 1.2060 1.2490 1.3516
   14.9 1.2052 1.2493 1.3519
   15.0 1.2044 1.2496 1.3523
   15.1 1.2037 1.2499 1.3526
   15.2 1.2029 1.2502 1.3530
   15.3 1.2021 1.2505 1.3533
   15.4 1.2013 1.2508 1.3537
   15.5 1.2005 1.2511 1.3540
   15.6 1.1998 1.2513 1.3543
   15.7 1.1990 1.2516 1.3547
   15.8 1.1982 1.2519 1.3550
   15.9 1.1974 1.2522 1.3553
   16.0 1.1967 1.2524 1.3556
   16.1 1.1959 1.2527 1.3559
   16.2 1.1952 1.2530 1.3562
   16.3 1.4051 1.2532 1.3565
   16.4 1.4041 1.2535 1.3568
   16.5 1.4030 1.2537 1.3571
   16.6 1.4019 1.2540 1.3574
   16.7 1.4009 1.2542 1.3577
   16.8 1.3998 1.2545 1.3580
   16.9 1.3988 1.2547 1.3583
   17.0 1.3977 1.2550 1.3586
   17.1 1.3967 1.2552 1.3589
   17.2 1.3956 1.2554 1.3591
   17.3 1.3946 1.2557 1.3594
   17.4 1.3935 1.2559 1.3597
   17.5 1.3925 1.2561 1.3599
   17.6 1.3915 1.2563 1.3602
   17.7 1.3904 1.2566 1.3605
   17.8 1.3894 1.2568 1.3607
   17.9 1.3884 1.2570 1.3610
   18.0 1.3874 1.2572 1.3612
   18.1 1.3863 1.2574 1.3615
   18.2 1.3853 1.2576 1.3617
   18.3 1.3843 1.2579 1.3620
   18.4 1.3833 1.2581 1.3622
   18.5 1.3823 1.2583 1.3624
   18.6 1.3813 1.2585 1.3627
   18.7 1.3803 1.2587 1.3629
   18.8 1.3793 1.2589 1.3631
   18.9 1.3783 1.2591 1.3634
   19.0 1.3773 1.2592 1.3636
   19.1 1.3764 1.2594 1.3638
   19.2 1.3754 1.2596 1.3640
   19.3 1.3744 1.2598 1.3643
   19.4 1.1556 1.2600 1.3645
   19.5 1.1552 1.2602 1.3647
   19.6 1.1548 1.2604 1.3649
   19.7 1.1544 1.2605 1.3651
   19.8 1.1540 1.2607 1.3653
   19.9 1.1536 1.2609 1.3655
   20.0 1.1532 1.2611 1.3657
   20.1 1.1528 1.2612 1.3659
   20.2 1.1524 1.2614 1.3661
   20.3 1.1520 1.2616 1.3663
   20.4 1.1516 1.2618 1.3665
   20.5 1.1512 1.2619 1.3667
   20.6 1.1508 1.2621 1.3669
   20.7 1.1504 1.2623 1.3671
   20.8 1.1500 1.2624 1.3673
   20.9 1.1496 1.2626 1.3675
   21.0 1.1492 1.2627 1.3677
   21.1 1.1488 1.2629 1.3679
   21.2 1.1484 1.2630 1.3681
   21.3 1.1480 1.2632 1.3682
   21.4 1.1476 1.2634 1.3684
   21.5 1.1473 1.2635 1.3686
   21.6 1.1469 1.2637 1.3688
   21.7 1.1465 1.2638 1.3689
   21.8 1.1461 1.2640 1.3691
   21.9 1.1457 1.2641 1.3693
   22.0 1.1453 1.2642 1.3695
   22.1 1.1449 1.2644 1.3696
   22.2 1.1446 1.2645 1.3698
   22.3 1.1442 1.2647 1.3700
   22.4 1.1438 1.2648 1.3701
   22.5 1.1434 1.2650 1.3703
   22.6 1.1430 1.2651 1.3705
   22.7 1.1427 1.2652 1.3706
   22.8 1.1423 1.2654 1.3708
   22.9 1.1419 1.2655 1.3709
   23.0 1.1415 1.2656 1.3711
   23.1 1.1412 1.2658 1.3712
   23.2 1.1408 1.2659 1.3714
   23.3 1.1404 1.2660 1.3715
   23.4 1.1401 1.2661 1.3717
   23.5 1.1397 1.2663 1.3718
   23.6 1.1393 1.2664 1.3720
   23.7 1.1389 1.2665 1.3721
   23.8 1.1386 1.2667 1.3723
   23.9 1.1382 1.2668 1.3724
   24.0 1.1378 1.2669 1.3726
}\gainctx
\addplot[ddrD, line width=0.8pt] table[x=S, y=integer] {\gainctx};
\addlegendentry{whole residents and riders}
\addplot[hbmD, line width=1.2pt, densely dashed] table[x=S, y=smooth] {\gainctx};
\addlegendentry{envelope (fractional $B$, $r$)}
\addplot[black!45, line width=0.6pt, densely dotted] table[x=S, y=smoothnc] {\gainctx};
\addlegendentry{envelope without HBM contention}
\addplot[only marks, mark=*, mark size=1.8pt, draw=black, fill=ddrD]
  coordinates {(14.724,1.207)};
\node[font=\sffamily\scriptsize, anchor=north] at (axis cs:14.724,1.195) {$B{=}3$, $r{=}2$};
\end{axis}
\end{tikzpicture}
\caption{\textbf{Gain against context} for Qwen2.5-72B (fp8 weights) on a
GH200, from the capacity model with $c$ and $T_0$ as functions of context
(\cref{eq:cS}), including the bandwidth the riders' copies take from HBM.
Residents and riders are whole requests, so the gain is a sawtooth around a
slowly rising envelope; the dotted line is the envelope if the copies cost HBM
nothing. Projection, not measurement.}
\label{fig:gain-ctx}
\end{figure}


Because residents and riders are whole requests, the gain does not vary
smoothly with context (\cref{fig:gain-ctx}). $B$ drops by a whole resident at
each context where one fewer fits, and the gain jumps with it. Riders change
more gently. As the context grows, the knee falls as $1/S$ and the last
rider's transfer is exposed a little more; the rider is kept while it still
pays, so a small rise in context or in measured step time costs a fraction of
a rider's worth, not a whole rider.

For Qwen2.5-72B (fp8 weights, $L{=}80$) on a GH200 the envelope rises slowly,
from $1.31\times$ at 8k tokens to $1.35\times$ at 24k, because the riders'
mechanism cost falls with context (\S\ref{sec:context}). The isolated
peaks above it, up to $1.58\times$ at 16.3k, lie where HBM strands most of one
resident's KV, so the baseline is at its weakest; we report the envelope and
treat those peaks as placement, not gain. At context 14\,724, for example,
$B = 3$ residents fit and $T_0 = 31.0$\,ms, so a token is worth
$T_0/B = 10.3$\,ms. Each rider costs
$c = 3.93$\,ms: 1.88\,ms of its own attention, 1.28\,ms of HBM bandwidth taken
by its copy, 0.14\,ms for its exposed first layer, and 0.63\,ms of mechanism.
The gate holds by $2.6\times$ and break-even is $\rmin = 1$. The link moves
one rider's KV in 11.5\,ms, so with $\delta = 6.8$\,ms the knee is
$\rstar = 3.75$: three riders hide, and a fourth exposes 1.9\,ms of its
transfer but still adds a token for less than the step time per token. Four
riders give $1.37\times$ throughput, against $1.20\times$ if only the riders
that hide behind $T_0$ alone were admitted; the envelope at this context is
$1.34\times$. What binds is the knee, not the gate. If the copies cost
nothing, the same point would give $1.45\times$ (dotted in
\cref{fig:gain-ctx}); \S\ref{sec:eval-proj} gives the projection's
sensitivity to each input.

\paragraph{What the link's bandwidth is worth.} The link adds $\Xlink/\mathit{BW}
= 10.5\%$ to the GPU's memory bandwidth (419\,GB/s on 4\,TB/s). A useful
measure of how well a design spends it is the throughput gained per unit of
bandwidth added,
\begin{equation}
  \eta \;=\; \frac{G - 1}{\Xlink/\mathit{BW}} .
  \label{eq:eta}
\end{equation}
Using the link only to read the \emph{same} requests' KV faster is bounded by
$\eta \le 1$: a step bound by memory bandwidth shortens by at most
$(\mathit{BW}+\Xlink)/\mathit{BW}$, before any contention. That is the regime
of bandwidth aggregation: DAK splits each operation's bytes between the tiers
and defines the bandwidth available as that of HBM plus the link~\cite{dak}.
DirectKV reads all KV over the link, and its efficiency is, in its own words,
``fundamentally limited by the bandwidth of the CPU--GPU
interconnect''~\cite{directkv}; when HBM suffices, the all-HBM engine is
faster~\cite{directkv}. A design that uses the link to add \emph{requests} can
exceed the bound, and riders do so.
% src: dak footnote 1: "Total available system memory bandwidth is defined by
%      GPU_HBM_BW + MIN(NVLink_C2C_BW, HOST_DRAM_BW)"; s4.2.1 optimal ratio
%      Bh/(Bh+Bg). directkv s7.4.3: "The efficiency of zero-copy offloading is
%      fundamentally limited by the bandwidth of the CPU-GPU interconnect";
%      s7.3.1: "SGLang provides the best performance when sufficient HBM is
%      available". pie s3.4.2 adaptive expansion grows the batch (quoted in s2).
Riders spend the link on \emph{more requests}, each of
which shares the weight read the batch already pays for; a rider's cost is its
own attention and its bytes, never a share of $W$. At the example above
$\eta = 3.5$, after paying the contention, and the envelope gives
$\eta = 2.9$--$3.4$ across 8k--24k tokens. The ratio of
\cref{eq:ratio} is largest where HBM is scarcest relative to the link, with
heavy weights and few residents. That is the regime the design exists for.
% src: paper/model/gain_curve.py (M22-corrected: c_m = 0.50*L/64, BW_kv = 0.64
%      HBM peak, BW_w = 0.75 peak, F = 3.39*L/64 = 4.24, X = 419 GB/s, b_tok
%      320 KiB, KAPPA 1.06 from E1, job 2191314): smooth 1.213 @8192, 1.249
%      @14724, 1.267 @24000; integer 1.405 @16300 (cap 2.98 -> B 2); @14724 B 3,
%      t0 30.99, c 3.932 = 0.625 + 1.885 + 1.279 + 0.144, r 2, G 1.199, eta
%      1.90; smooth 1.203 @8192, 1.238 @14724, 1.256 @24000; integer 1.395
%      @16300; eta_smooth 1.94 @8192, 2.44 @24000. Copies free: G 1.283.
%      Replaces 1.28x (2026-10-20), which replaced 1.37x.
%      2026-10-22, past-knee rule (rule="max", delta 6.8): @14724 r 4, G 1.368,
%      eta 3.51, exposed 1.9 ms, r* 3.75; r 3 1.318, r 5 1.284; envelope 1.308
%      @8192, 1.337 @14724, 1.351 @24000 (eta 2.94..3.35); integer 1.583
%      @16300; copies free (kappa 0, no first layer) 1.450 (r 3).
On a GH200, $T_0$ would be measured, not computed. Here it comes from
\cref{eq:cS} with the bandwidth efficiencies measured on A100; the projection's
sensitivity to it, and to every other input, is in \S\ref{sec:eval-proj}.
